\documentclass[12pt]{article}
\usepackage[utf8]{inputenc}
\usepackage[T1]{fontenc}
\usepackage[margin=1in]{geometry}
\usepackage{setspace}
\usepackage{xcolor}
\usepackage[colorlinks=true, urlcolor=blue!60!black, linkcolor=blue!60!black]{hyperref}
\setlength{\parindent}{2em}
\setlength{\parskip}{1.25em}
\renewcommand{\baselinestretch}{1.0}
\begin{document}
\begin{center}
{\LARGE \textit{Obergefell v. Hodges}, 576 U.S. 644 (2015)\par}
\vspace{0.8em}
{\large Case Brief\par}
\end{center}
\vspace{1em}

\section{Memory Jogger}

Rejecting a rigid history-and-tradition ceiling for substantive due process, the Supreme Court held 5--4 that the interlocking guarantees of the Fourteenth Amendment's Due Process and Equal Protection Clauses require every state to license same-sex marriages and recognize same-sex marriages lawfully performed out of state (\href{https://www.courtlistener.com/opinion/2812209/obergefell-v-hodges/}{CourtListener}).

\section{Detailed Case Facts}

Fourteen same-sex couples and two men whose same-sex partners had died challenged state constitutional amendments and statutes in Michigan, Kentucky, Ohio, and Tennessee that defined civil marriage exclusively as a union between one man and one woman and refused to recognize same-sex marriages lawfully performed in other jurisdictions (\href{https://www.law.cornell.edu/supremecourt/text/14-556}{Cornell LII}). Specifically, the challenged provisions included Mich. Const., Art. I, \S~25; Ky. Const. \S~233A; Ohio Rev. Code Ann. \S~3101.01; and Tenn. Const., Art. XI, \S~18 (\href{https://www.law.cornell.edu/supremecourt/text/14-556}{Cornell LII}).

Justice Kennedy's majority opinion highlighted three petitioners' stories to illustrate the concrete legal and dignitary harms inflicted by the state bans (\href{https://www.law.cornell.edu/supremecourt/text/14-556}{Cornell LII}):
\begin{itemize}
    \item \textbf{James Obergefell and John Arthur (Ohio):} Together for over two decades, Obergefell and Arthur resolved to marry after Arthur was diagnosed with amyotrophic lateral sclerosis (ALS) in 2011. Because Ohio barred same-sex marriage and Arthur could barely move, they traveled in a medical transport plane to Baltimore, Maryland, where same-sex marriage was legal, and wed on the tarmac in July 2013. Three months later, Arthur died, and Ohio law barred Obergefell from being listed as the surviving spouse on Arthur's death certificate, leaving them legally ``strangers even in death'' (\href{https://www.law.cornell.edu/supremecourt/text/14-556}{Cornell LII}).
    \item \textbf{April DeBoer and Jayne Rowse (Michigan):} Two nurses who held a commitment ceremony in 2007 fostered and adopted three special-needs children. Because Michigan allowed only opposite-sex married couples or single individuals to adopt, each child could have only one woman as a legal parent, leaving the family vulnerable in medical emergencies or if either partner died (\href{https://www.law.cornell.edu/supremecourt/text/14-556}{Cornell LII}).
    \item \textbf{Ijpe DeKoe and Thomas Kostura (Tennessee):} Army Reserve Sergeant First Class DeKoe married Kostura in New York in 2011 shortly before deploying to Afghanistan for nearly a year. Upon returning and settling in Tennessee, where DeKoe served full-time in the Army Reserve, Tennessee stripped their lawful marriage of legal recognition whenever they resided within state borders (\href{https://www.law.cornell.edu/supremecourt/text/14-556}{Cornell LII}).
\end{itemize}
Additional plaintiffs in Kentucky (\textit{Bourke v. Beshear} and \textit{Love v. Beshear}), Ohio (\textit{Henry v. Himes}), and Tennessee (\textit{Tanco v. Haslam}) challenged analogous licensing and out-of-state recognition bans (\href{https://www.theusconstitution.org/litigation/de-boer-v-snyder-6th-cir/}{Constitutional Accountability Center}).

\section{Procedural History}

The plaintiffs filed six separate federal lawsuits across four states in U.S. District Courts: \textit{Obergefell v. Wymyslo}, 962 F. Supp. 2d 968 (S.D. Ohio 2013), and \textit{Henry v. Himes}, 14 F. Supp. 3d 1036 (S.D. Ohio 2014); \textit{DeBoer v. Snyder}, 973 F. Supp. 2d 757 (E.D. Mich. 2014); \textit{Bourke v. Beshear}, 996 F. Supp. 2d 542 (W.D. Ky. 2014), and \textit{Love v. Beshear}, 989 F. Supp. 2d 536 (W.D. Ky. 2014); and \textit{Tanco v. Haslam}, 7 F. Supp. 3d 759 (M.D. Tenn. 2014) (\href{https://www.law.cornell.edu/supremecourt/text/14-556}{Cornell LII}). Every district court ruled in favor of the plaintiffs, holding that the state bans violated the Fourteenth Amendment (\href{https://www.oyez.org/cases/2014/14-556}{Oyez}).

State officials appealed to the U.S. Court of Appeals for the Sixth Circuit, which consolidated the cases under \textit{DeBoer v. Snyder}, 772 F.3d 388 (6th Cir. 2014) (\href{https://www.law.cornell.edu/supremecourt/text/14-556}{Cornell LII}). On November 6, 2014, a 2--1 panel---Judge Jeffrey Sutton writing for himself and Judge Deborah Cook, over Judge Martha Daughtrey's dissent---reversed all six district courts (\href{https://www.theusconstitution.org/litigation/de-boer-v-snyder-6th-cir/}{Constitutional Accountability Center}). The Sixth Circuit concluded it was bound by the Supreme Court's 1972 one-line summary dismissal in \textit{Baker v. Nelson}, 409 U.S. 810 (1972), and held that the Fourteenth Amendment left the definition of marriage to state voters (\href{https://www.theusconstitution.org/litigation/de-boer-v-snyder-6th-cir/}{Constitutional Accountability Center}). Because the Fourth, Seventh, Ninth, and Tenth Circuits had struck down similar bans, the Sixth Circuit created a direct circuit split (\href{https://www.theusconstitution.org/litigation/de-boer-v-snyder-6th-cir/}{Constitutional Accountability Center}).

The Supreme Court granted certiorari on January 16, 2015, consolidating four petitions---Nos. 14--556 (\textit{Obergefell}), 14--562 (\textit{Tanco}), 14--571 (\textit{DeBoer}), and 14--574 (\textit{Bourke})---limited to two questions: (1) whether the Fourteenth Amendment requires a state to license a marriage between two people of the same sex, and (2) whether it requires a state to recognize a same-sex marriage lawfully licensed and performed out of state (\href{https://supreme.justia.com/cases/federal/us/576/644/}{Justia}). Oral argument took place on April 28, 2015, and the Court issued its decision on June 26, 2015 (\href{https://supreme.justia.com/cases/federal/us/576/644/}{Justia}).

\section{Judicial Votes}

The Court ruled \textbf{5--4} to reverse the Sixth Circuit (\href{https://www.scotusblog.com/cases/obergefell-v-hodges/}{SCOTUSblog}):
\begin{itemize}
    \item \textbf{Majority (5):} Justice Anthony M. Kennedy authored the opinion of the Court, joined in full by Justices Ruth Bader Ginsburg, Stephen G. Breyer, Sonia Sotomayor, and Elena Kagan (\href{https://www.law.cornell.edu/supremecourt/text/14-556}{Cornell LII}). None of the four joining Justices wrote a separate concurrence.
    \item \textbf{Dissents (4):} Every dissenting Justice wrote separately, with overlapping joins (\href{https://www.scotusblog.com/cases/obergefell-v-hodges/}{SCOTUSblog}):
    \begin{itemize}
        \item Chief Justice John G. Roberts, Jr. filed a dissenting opinion, joined by Justices Antonin Scalia and Clarence Thomas.
        \item Justice Antonin Scalia filed a dissenting opinion, joined by Justice Clarence Thomas.
        \item Justice Clarence Thomas filed a dissenting opinion, joined by Justice Antonin Scalia.
        \item Justice Samuel A. Alito, Jr. filed a dissenting opinion, joined by Justices Antonin Scalia and Clarence Thomas.
    \end{itemize}
\end{itemize}

\section{Holding}

\textbf{Judgment:} Reversed (\href{https://www.law.cornell.edu/supremecourt/text/14-556}{Cornell LII}). \textit{Baker v. Nelson}, 409 U.S. 810 (1972), is expressly overruled (\href{https://www.law.cornell.edu/supremecourt/text/14-556}{Cornell LII}).

\textbf{Rule Announced:} Under both the Due Process Clause and the Equal Protection Clause of the Fourteenth Amendment, the right to marry is a fundamental right inherent in the liberty of the person, and same-sex couples may not be deprived of that right (\href{https://www.law.cornell.edu/supremecourt/text/14-556}{Cornell LII}). Accordingly, states must (1) license marriages between two people of the same sex on the same terms and conditions as opposite-sex couples, and (2) recognize lawful same-sex marriages licensed and performed in other states (\href{https://www.law.cornell.edu/supremecourt/text/14-556}{Cornell LII}). Methodologically, in identifying unenumerated fundamental rights under the Due Process Clause, courts exercise ``reasoned judgment'' in which history and tradition guide and discipline the inquiry but do not set its outer boundaries, rather than applying the narrow ``careful description'' formula of \textit{Washington v. Glucksberg}, 521 U.S. 702 (1997), to marriage and intimate association (\href{https://www.law.cornell.edu/supremecourt/text/14-556}{Cornell LII}).

\textbf{Holding vs. Dicta:} The binding holding invalidates state bans on licensing and recognizing civil marriage for two consenting adults of the same sex, mandates equal access to the state-conferred constellation of marital benefits (later enforced in \textit{Pavan v. Smith}, 582 U.S. 563 (2017)), and overrules \textit{Baker v. Nelson} (\href{https://www.law.cornell.edu/supremecourt/text/14-556}{Cornell LII}; \href{https://www.americanbar.org/groups/family_law/resources/family-law-quarterly/2025-december/obergefell-victories-setbacks-protecting-rights-lgbtq-families/}{ABA Family Law Quarterly}). By contrast, several prominent passages are non-binding dicta or unresolved questions: (1) Justice Kennedy's assurances regarding First Amendment protections for religious organizations and believers who ``advocate'' and ``teach'' traditional views of marriage did not adjudicate any concrete Free Exercise or compelled-speech dispute; (2) the Court's observation that the cases involved ``only the rights of two consenting adults'' did not formally adjudicate plural marriage; and (3) the Court declined to decide whether sexual orientation constitutes a suspect or quasi-suspect classification triggering heightened (intermediate or strict) scrutiny under traditional tiered Equal Protection doctrine, relying instead on the ``synergy'' between Due Process and Equal Protection (\href{https://www.law.cornell.edu/supremecourt/text/14-556}{Cornell LII}).

\section{Analysis of Opinions}

\subsection{Majority Opinion (Justice Kennedy)}

Justice Kennedy framed the case around the evolution of marriage and the proper methodology for substantive due process (\href{https://www.law.cornell.edu/supremecourt/text/14-556}{Cornell LII}). Tracing marriage from arranged unions and the doctrine of coverture to modern egalitarian partnerships, he argued that historical continuity coexists with transformative change that strengthens the institution (\href{https://www.law.cornell.edu/supremecourt/text/14-556}{Cornell LII}).

On substantive due process methodology, Kennedy openly rejected respondents' insistence that \textit{Washington v. Glucksberg}, 521 U.S. 702 (1997), required defining the asserted right narrowly as a ``right to same-sex marriage'' and testing whether it was deeply rooted in 1868 history (\href{https://www.law.cornell.edu/supremecourt/text/14-556}{Cornell LII}). Invoking Justice Harlan's dissent in \textit{Poe v. Ullman}, 367 U.S. 497, 542 (1961), Kennedy maintained that identifying fundamental rights ``has not been reduced to any formula'' and requires ``reasoned judgment'' (\href{https://www.law.cornell.edu/supremecourt/text/14-556}{Cornell LII}). Just as \textit{Loving v. Virginia}, 388 U.S. 1 (1967), \textit{Zablocki v. Redhail}, 434 U.S. 374 (1978), and \textit{Turner v. Safley}, 482 U.S. 78 (1987), asked about the ``right to marry in its comprehensive sense'' rather than a right to interracial, deadbeat-parent, or prisoner marriage, \textit{Glucksberg}'s narrow framing could not allow past exclusions to serve as their own permanent justification (\href{https://www.law.cornell.edu/supremecourt/text/14-556}{Cornell LII}).

Kennedy then articulated four principles demonstrating why marriage is fundamental and why each applies with equal force to same-sex couples (\href{https://www.law.cornell.edu/supremecourt/text/14-556}{Cornell LII}):
\begin{enumerate}
    \item \textbf{Individual Autonomy:} Personal choice regarding marriage is among the most intimate acts of self-definition and autonomy.
    \item \textbf{Two-Person Intimate Union:} Marriage supports a unique, committed two-person bond unlike any other, extending liberty beyond mere decriminalization under \textit{Lawrence v. Texas}, 539 U.S. 558 (2003).
    \item \textbf{Safeguarding Children and Families:} Marriage confers legal stability and social recognition on children raised by same-sex couples, sparing them material burdens and the stigma of being deemed ``lesser,'' while not conditioning marriage on procreative capacity.
    \item \textbf{Keystone of Social Order:} States have attached a vast constellation of rights, benefits, and responsibilities to marriage (taxation, inheritance, hospital access, adoption, death certificates), making exclusion demeaning and destabilizing.
\end{enumerate}

Turning to Equal Protection, Kennedy avoided formal tiered-scrutiny analysis (neither declaring sexual orientation a suspect/quasi-suspect class nor applying standard rational-basis review) and instead emphasized the ``interlocking nature'' and ``synergy'' of the Due Process and Equal Protection Clauses (\href{https://www.law.cornell.edu/supremecourt/text/14-556}{Cornell LII}). Citing \textit{Loving}, \textit{Zablocki}, \textit{M.L.B. v. S.L.J.}, 519 U.S. 102 (1996), and \textit{Lawrence}, he concluded that denying same-sex couples a fundamental right both burdens liberty and works a grave, subordinating inequality (\href{https://www.law.cornell.edu/supremecourt/text/14-556}{Cornell LII}). Finally, Kennedy rejected arguments to await further democratic deliberation, quoting \textit{West Virginia Board of Education v. Barnette}, 319 U.S. 624, 638 (1943), that ``fundamental rights may not be submitted to vote,'' dismissed claims that same-sex marriage would harm opposite-sex marriage, reaffirmed First Amendment protections for religious teaching, and held that recognition of out-of-state same-sex marriages necessarily followed (\href{https://www.law.cornell.edu/supremecourt/text/14-556}{Cornell LII}).

\subsection{Dissenting Opinion (Chief Justice Roberts)}

Chief Justice Roberts, joined by Justices Scalia and Thomas, cast the case not as a policy dispute over marriage, but as a structural question of ``who decides'' in a democratic republic (\href{https://www.law.cornell.edu/supremecourt/text/14-556}{Cornell LII}). Invoking \textit{Dred Scott v. Sandford}, 19 How. 393 (1857), and \textit{Lochner v. New York}, 198 U.S. 45 (1905), Roberts accused the majority of reviving unprincipled substantive due process and ``effectively overrul[ing] \textit{Glucksberg}'' to enact its own moral philosophy---quoting Judge Henry Friendly that the Fourteenth Amendment no more enacts John Stuart Mill's \textit{On Liberty} than Herbert Spencer's \textit{Social Statics} (\href{https://www.law.cornell.edu/supremecourt/text/14-556}{Cornell LII}).

Roberts distinguished \textit{Loving}, \textit{Zablocki}, and \textit{Turner} as removing barriers to marriage as traditionally defined, whereas petitioners sought to compel states to redefine the core opposite-sex institution itself (\href{https://www.law.cornell.edu/supremecourt/text/14-556}{Cornell LII}). He argued that the majority's autonomy and anti-stigma rationale offered no principled boundary against plural marriage, criticized the Equal Protection ``synergy'' section as devoid of standard means-ends scrutiny, and warned that short-circuiting democratic debate would provoke backlash and jeopardize the Free Exercise rights of religious institutions---noting pointedly that the majority protected the right to ``advocate'' and ``teach,'' but omitted the Constitution's word ``exercise'' (\href{https://www.law.cornell.edu/supremecourt/text/14-556}{Cornell LII}).

\subsection{Dissenting Opinion (Justice Scalia)}

Justice Scalia, joined by Justice Thomas, joined Chief Justice Roberts in full and wrote separately to denounce the decision as a ``judicial Putsch'' and a ``threat to American democracy'' (\href{https://www.law.cornell.edu/supremecourt/text/14-556}{Cornell LII}). Applying original public meaning, Scalia contended that because every state in 1868 defined marriage as between one man and one woman without constitutional controversy, the Fourteenth Amendment cannot plausibly prohibit that definition (\href{https://www.law.cornell.edu/supremecourt/text/14-556}{Cornell LII}). He attacked the unrepresentative demographics of the nine-member Court---all Harvard or Yale Law graduates, predominantly from coastal states---arguing that imposing social transformation by judicial decree violates ``no social transformation without representation,'' and mocked the majority's prose as ``the mystical aphorisms of the fortune cookie'' (\href{https://www.law.cornell.edu/supremecourt/text/14-556}{Cornell LII}).

\subsection{Dissenting Opinion (Justice Thomas)}

Justice Thomas, joined by Justice Scalia, rejected substantive due process altogether as a ``dangerous fiction'' untethered from constitutional text (\href{https://www.law.cornell.edu/supremecourt/text/14-556}{Cornell LII}). Tracing ``liberty'' from Magna Carta through Blackstone, Coke, and early state constitutions, Thomas argued that the Due Process Clauses protect only negative liberty---specifically, freedom from physical restraint or, at most, freedom from governmental interference---never a positive entitlement to governmental recognition, licenses, or benefits (\href{https://www.law.cornell.edu/supremecourt/text/14-556}{Cornell LII}). He also warned of severe collateral damage to religious liberty and rejected the majority's concept of state-conferred dignity, insisting under the Declaration of Independence that human dignity is innate and ``cannot be taken away by the government'' (\href{https://www.law.cornell.edu/supremecourt/text/14-556}{Cornell LII}).

\subsection{Dissenting Opinion (Justice Alito)}

Justice Alito, joined by Justices Scalia and Thomas, anchored his dissent in \textit{Glucksberg}'s ``deeply rooted in this Nation's history and tradition'' standard, noting that no country permitted same-sex marriage before the Netherlands in 2000 and no U.S. state did so before Massachusetts in 2003 (\href{https://www.law.cornell.edu/supremecourt/text/14-556}{Cornell LII}). He argued that states had a rational secular basis for adhering to the traditional, procreation-linked understanding of marriage rather than the majority's ``postmodern,'' fulfillment-centered model, and cautioned that the ruling would be used to ``vilify'' traditional believers as bigots and erode federalism's capacity to let citizens with differing views coexist (\href{https://www.law.cornell.edu/supremecourt/text/14-556}{Cornell LII}).

\section{Examples: Future Applications}

\begin{enumerate}
    \item \textbf{Same-Side (Invalidated under \textit{Obergefell}): State Refusal to List Both Same-Sex Spouses on a Child's Birth Certificate.} A state statute automatically lists a birth mother's opposite-sex husband as the legal father on a child's birth certificate even when the child was conceived via anonymous donor insemination, but refuses to list a birth mother's same-sex wife without a costly adoption proceeding. Under \textit{Obergefell}, because states must afford same-sex couples civil marriage ``on the same terms and conditions'' and with the same ``constellation of benefits'' linked to marital status---explicitly including birth certificates---this exclusion is unconstitutional (as the Court summarily confirmed in \textit{Pavan v. Smith}, 582 U.S. 563 (2017)) (\href{https://www.law.cornell.edu/supremecourt/text/14-556}{Cornell LII}; \href{https://www.americanbar.org/groups/family_law/resources/family-law-quarterly/2025-december/obergefell-victories-setbacks-protecting-rights-lgbtq-families/}{ABA Family Law Quarterly}).

    \item \textbf{Same-Side (Invalidated under \textit{Obergefell}): Municipal Denial of Spousal Health and Survivor Pension Benefits to Same-Sex Married Employees.} A city extends taxpayer-funded health insurance and survivor pension benefits to opposite-sex spouses of municipal firefighters, but denies those benefits to same-sex spouses lawfully married in another state on the ground that local ordinances predate 2015. Because \textit{Obergefell} holds that states and their political subdivisions must recognize lawful same-sex marriages and cannot withhold the material incidents of marriage (specifically enumerating health insurance and survivor benefits) on the basis of a marriage's same-sex character, the denial violates the Fourteenth Amendment (\href{https://www.law.cornell.edu/supremecourt/text/14-556}{Cornell LII}).

    \item \textbf{Opposite-Side (Upheld under \textit{Obergefell}): Criminal Ban on Physician-Assisted Suicide.} A terminally ill adult challenges a state statute criminalizing physician-assisted suicide, arguing under \textit{Obergefell}'s ``reasoned judgment'' and individual-autonomy framework that deciding the timing and manner of death is central to personal dignity. Even after \textit{Obergefell}, this claim fails because Justice Kennedy expressly distinguished \textit{Washington v. Glucksberg}, 521 U.S. 702 (1997), stating that \textit{Glucksberg}'s narrow historical approach ``may have been appropriate for the asserted right there involved (physician-assisted suicide),'' and the Court later reaffirmed \textit{Glucksberg}'s history-and-tradition test for non-marriage substantive due process claims in \textit{Dobbs v. Jackson Women's Health Organization}, 597 U.S. 215 (2022) (\href{https://www.law.cornell.edu/supremecourt/text/14-556}{Cornell LII}; \href{https://www.americanbar.org/groups/family_law/resources/family-law-quarterly/2025-december/obergefell-victories-setbacks-protecting-rights-lgbtq-families/}{ABA Family Law Quarterly}).

    \item \textbf{Opposite-Side (Upheld under \textit{Obergefell}): Minimum Age and Close-Consanguinity Marriage Restrictions.} A 15-year-old or an adult seeking to marry a biological sibling challenges state marriage restrictions under \textit{Obergefell}'s autonomy principle. Both challenges fail because \textit{Obergefell} explicitly limited its reasoning to ``two consenting adults whose marriages would pose no risk of harm to themselves or third parties,'' preserving states' compelling interests in protecting minors from immature consent and preventing the familial and genetic harms historically associated with incestuous unions (\href{https://www.law.cornell.edu/supremecourt/text/14-556}{Cornell LII}).

    \item \textbf{Fence-Sitter (Close / Unresolved under \textit{Obergefell}): Constitutional Claim to Civil Plural (Polyamorous) Marriage.} Three consenting adults living in a committed, long-term relationship and co-parenting children challenge a state statute limiting civil marriage to two persons. On one side, Chief Justice Roberts's dissent pointed out that \textit{Obergefell}'s autonomy, anti-stigma, and child-welfare principles could be invoked by plural families, and scholars have debated whether the ``two-person'' premise is merely a historical artifact (\href{https://www.law.cornell.edu/supremecourt/text/14-556}{Cornell LII}; \href{https://newenglrev.com/2015/08/31/obergefell-and-the-future-of-plural-marriage/}{New England Law Review}). On the other side, Justice Kennedy repeatedly anchored his second principle in a ``two-person union'' and noted that plural marriage raises distinct administrative, third-party, and equity complexities (such as custody, divorce, inheritance, and power asymmetries) that states may regulate, leaving the issue genuinely contested under \textit{Obergefell}'s logic (\href{https://www.law.cornell.edu/supremecourt/text/14-556}{Cornell LII}).
\end{enumerate}

\section{Critique}

\textbf{Originalist and Textualist Perspectives:} Originalist and textualist scholars largely echo the four dissents in contending that \textit{Obergefell} lacks grounding in the original public meaning of the Fourteenth Amendment. Critics argue that in 1868, neither the Due Process Clause nor the Equal Protection Clause was understood to require states to license same-sex marriage, and that Kennedy's ``reasoned judgment'' methodology displaces \textit{Glucksberg}'s objective historical constraints with judicial value judgments (\href{https://dartmouthlawjournal.org/article/18840-obergefell-v-hodges-rethinking-justice-scalia-s-originalism.pdf}{Dartmouth Law Journal}; \href{https://fedsoc.org/case/obergefell-v-hodges}{Federalist Society}). Following \textit{Dobbs v. Jackson Women's Health Organization}, 597 U.S. 215 (2022)---which revived \textit{Glucksberg}'s ``deeply rooted in this Nation's history and tradition'' framework---Justice Thomas explicitly urged the Court to ``reconsider all of this Court's substantive due process precedents, including \textit{Griswold}, \textit{Lawrence}, and \textit{Obergefell},'' prompting Congress to enact the Respect for Marriage Act of 2022 to provide a federal statutory backstop (\href{https://www.americanbar.org/groups/family_law/resources/family-law-quarterly/2025-december/obergefell-victories-setbacks-protecting-rights-lgbtq-families/}{ABA Family Law Quarterly}; \href{https://community.lawschool.cornell.edu/wp-content/uploads/2024/08/Scaglia-final.pdf}{Cornell Law School Community}).

\textbf{Progressive and Doctrinal Perspectives:} Many progressive constitutional scholars celebrated the result as a landmark triumph for LGBTQ equality while offering nuanced doctrinal critiques. Kenji Yoshino praised \textit{Obergefell} as a ``game-changer'' for substantive due process that replaced \textit{Glucksberg}'s rigid formula with Justice Harlan's common-law ``antisubordination'' liberty approach, while noting that the Court's refusal to articulate a formal tier of Equal Protection scrutiny for sexual orientation left lower courts without clear guidance in non-marriage discrimination cases (\href{https://harvardlawreview.org/print/vol-129/a-new-birth-of-freedom-obergefell-v-hodges/}{Harvard Law Review}; \href{https://www.scotusblog.com/2015/11/academic-highlight-yoshino-on-obergefell-v-hodges/}{SCOTUSblog}). Laurence Tribe similarly lauded the opinion's doctrinal fusion of liberty and equality into ``equal dignity,'' arguing that the two clauses form a ``double helix'' rather than siloed boxes (\href{https://harvardlawreview.org/wp-content/uploads/2015/11/vol129_Tribe.pdf}{Harvard Law Review Forum}). Yet other progressive and family-law scholars criticized Kennedy's rhetoric for glorifying civil marriage as the pinnacle of human fulfillment, warning that portraying unmarried individuals as ``condemned to live in loneliness'' demeans single parents and nonmarital families and undermines the opinion's own commitment to equal dignity (\href{https://www.cardozolawreview.com/whats-wrong-with-obergefell/}{Cardozo Law Review}).

\textbf{Underlying Competing Values:} At bottom, \textit{Obergefell} exposes an enduring clash between (1) individual autonomy, equal dignity, and countermajoritarian judicial protection of marginalized minorities, and (2) democratic self-governance, federalism, judicial restraint, and religious liberty. The methodological war between \textit{Poe}/\textit{Obergefell} (evolving reasoned judgment) and \textit{Glucksberg}/\textit{Dobbs} (strict historical tradition) reflects fundamentally competing visions of how the judiciary should mediate social change under the Fourteenth Amendment.

\section{Key Quotations}

\begin{quote}
``The identification and protection of fundamental rights is an enduring part of the judicial duty to interpret the Constitution. That responsibility, however, `has not been reduced to any formula.' \textit{Poe v. Ullman}, 367 U.S. 497, 542 (1961) (Harlan, J., dissenting). Rather, it requires courts to exercise reasoned judgment in identifying interests of the person so fundamental that the State must accord them its respect. [\dots] History and tradition guide and discipline this inquiry but do not set its outer boundaries. [\dots] That method respects our history and learns from it without allowing the past alone to rule the present.''
\end{quote}
\noindent --- Justice Kennedy, Opinion of the Court, 576 U.S. at 664 (\href{https://www.law.cornell.edu/supremecourt/text/14-556}{Cornell LII}; \href{https://www.courtlistener.com/opinion/2812209/obergefell-v-hodges/}{CourtListener}).

\begin{quote}
``If rights were defined by who exercised them in the past, then received practices could serve as their own continued justification and new groups could not invoke rights once denied.''
\end{quote}
\noindent --- Justice Kennedy, Opinion of the Court, 576 U.S. at 671 (\href{https://www.law.cornell.edu/supremecourt/text/14-556}{Cornell LII}).

\begin{quote}
``No union is more profound than marriage, for it embodies the highest ideals of love, fidelity, devotion, sacrifice, and family. In forming a marital union, two people become something greater than once they were. As some of the petitioners in these cases demonstrate, marriage embodies a love that may endure even past death. It would misunderstand these men and women to say they disrespect the idea of marriage. Their plea is that they do respect it, respect it so deeply that they seek to find its fulfillment for themselves. Their hope is not to be condemned to live in loneliness, excluded from one of civilization's oldest institutions. They ask for equal dignity in the eyes of the law. The Constitution grants them that right.''
\end{quote}
\noindent --- Justice Kennedy, Opinion of the Court, 576 U.S. at 681 (\href{https://www.law.cornell.edu/supremecourt/text/14-556}{Cornell LII}; \href{https://supreme.justia.com/cases/federal/us/576/644/}{Justia}).

\begin{quote}
``If you are among the many Americans---of whatever sexual orientation---who favor expanding same-sex marriage, by all means celebrate today's decision. Celebrate the achievement of a desired goal. Celebrate the opportunity for a new expression of commitment to a partner. Celebrate the availability of new benefits. But do not celebrate the Constitution. It had nothing to do with it.''
\end{quote}
\noindent --- Chief Justice Roberts, dissenting, 576 U.S. at 713 (\href{https://www.law.cornell.edu/supremecourt/text/14-556}{Cornell LII}).

\begin{quote}
``Today's decree says that my Ruler, and the Ruler of 320 million Americans coast-to-coast, is a majority of the nine lawyers on the Supreme Court. The opinion in these cases is the furthest extension in fact---and the furthest extension one can even imagine---of the Court's claimed power to create `liberties' that the Constitution and its Amendments neglect to mention.''
\end{quote}
\noindent --- Justice Scalia, dissenting, 576 U.S. at 713 (\href{https://www.law.cornell.edu/supremecourt/text/14-556}{Cornell LII}).

\end{document}
